\documentclass[10pt]{article}
\usepackage{icmlko}

\icmlmeta{IP 파운데이션 월드모델}{482}{자와 판이 다른 시대를 산다}{크롤 229시간을 태우기 전에 기울기를 쟀다 --- 그런데 자를 판에 대 보니 3행이 겹쳤다. 그리고 「텍스트가 있다」는 이미 판의 열이었다}

\begin{document}

\twocolumn[
\icmlheader

\begin{abstract}
판 $\rho$ 가 400사이클 넘게 \textbf{0.46--0.47} 에서 평평한데 다음 블록의 계획은
\textbf{광역 크롤}이다(만화 전량 12.5h · 12도메인 34.4h · 댓글까지 229h). 직전 세
노트는 전부 \emph{``텍스트를 어디서 어떻게 긁나''} 를 쟀고 \emph{``텍스트가 들어오면
$\rho$ 가 오르나''} 는 한 번도 안 쟀다. 그래서 \textbf{크롤 전에 기울기를} 쟀다.

\textbf{0단계가 실험을 두 번 바꿨다.} ① 직전 노트의 자(만화 롱테일 220건)와 판의
유보(3,775행 · 12도메인)의 \textbf{교집합은 3행}이다 --- 자는 전 연도 롱테일에서
왔고 판은 \textbf{2025년 이후만} 채점한다. 자와 판이 다른 시대를 산다.
② 요청 0 으로 유보에 닿는 용량을 다시 짓다가, \textbf{「문서가 있다」는 이미 판의
열}이라는 것을 \emph{측정 전에} 알았다 --- 판은 축마다 값 열과 \textbf{관측 표시자
열}을 함께 쓰고, 위키 축의 표시자가 곧 문서 존재다.

그래서 용량을 \textbf{둘로 갈라} 쟀다(존재 이진 · 양). 관찰(팔 A)에서
\textbf{존재는 죽고}($\rho_A=-0.0097$ · 순열 널 $2\sigma=0.0406$ · $p=0.607$)
\textbf{양만 산다}($+0.0692$ · BCa 95\% 하한 $+0.0020$ · 상한 $+0.1368$ ·
글자수·어절수 통제 후
$+0.0704$ 로 \textbf{유지율 1.02}). 그런데 개입(팔 B)에서는 \textbf{둘 다 문턱
0.00353 을 못 넘는다} --- 존재 $+0.00318$(문턱의 0.90배) · 양 $\mathbf{-0.00152}$,
그리고 양은 \textbf{세계애니 도메인에서 $-0.015$ 로 되레 해롭다}. 관찰에서 제일 센
도메인이 개입에서 제일 크게 손해를 본다.

\textbf{그래서 크롤은 지금 근거가 없다.} 다만 우리가 적을 말은 \emph{``텍스트는
쓸모없다''} 가 아니다 --- \textbf{이 자로는 이 판을 못 넘었다}, 그리고 자가 재는
것이 \textbf{텍스트가 아니라 텍스트의 대용값}이다.
\end{abstract}
]

\section{0단계 --- 자를 판에 대 보니 3행이었다}

조항 60 은 \emph{``두 수를 이어 붙일 때 분모가 같은지 확인하라''} 다. 직전 노트의
자는 \textbf{만화 220건}(실물 200 + 음성 20)에서 만들어졌고 판의 유보는
\textbf{3,775행 · 12도메인}이다. 세어 봤다(그림 \ref{fig:1} 왼쪽).

\begin{table}[h]\centering\small
\begin{tabular}{lr}
\hline
895 정답지 실물 & 200 \\
그중 \texttt{manga\_records}(7,982) 안 & 200 \\
그중 \textbf{판의 만화 행}(5,905) 안 & 141 \\
그중 \textbf{유보}(258 · $t\ge2025.0$) 안 & \textbf{3} \\
\hline
\end{tabular}
\end{table}

원인은 \textbf{시간축}이다. 판의 만화 유보는 2025.02--2026.49 시작작이고, 자의 200 은
1958--2026 롱테일 표본이다(2025년 이후가 12건, 그중 판 행에 있는 것이 3건).
\claimx{895 의 자와 판의 유보는 교집합이 3행이다 --- 자를 그대로 열로 주입하는
설계는 산술적으로 불가능하다. 스피어만 최소 행수는 20 이다.}{만화 유보 258 중 895 표본에 든 것을 다시 세어 20행이 넘으면 이 주장은 죽는다. 반대로 자를 유보에서 다시 뽑으면 교집합이 정의상 258 이 된다.}

이건 \emph{``기울기가 0''} 이 아니라 \textbf{``이 자로는 이 판을 못 잰다''} 다
(조항 59). 만화 하나로 판 문턱 0.00353 을 넘으려면 만화 도메인
$\Delta\rho \ge 0.00353\times3775/258 = \mathbf{0.0517}$ 이 필요하다.
\claimx{자를 지을 때 표본을 \textbf{유보에서} 뽑아야 한다. 895 는 롱테일 전
연도에서 뽑았고 판은 2025년 이후만 채점한다.}{유보에서 뽑은 표본으로 같은 자를 다시 세웠는데도 판 $\Delta\rho$ 가 문턱을 못 넘으면, 못 넘는 원인은 표본이 아니라 자다.}

\section{용량은 하나가 아니다 --- 그리고 하나는 이미 판 안에 있었다}

직전 사이클의 티처가 확정한 것: \emph{양(바이트)은 아무것도 안 더한다. 존재 이진이
바이트보다 낫다}(0.5365 대 0.5291). 그래서 이번엔 \textbf{존재와 양을 갈라} 쟀다.

요청 0 으로 유보에 닿는 것을 전수 조사했다. \textbf{용량 ㄱ 존재 이진}은 위키 문서
캐시(\texttt{page}$\ne$\texttt{null})에서 7도메인 2,502행, \textbf{용량 ㄴ 양}은
설명문 글자수(세계애니 300 + 펀딩 529)다. 만화는 \textbf{둘 다 못 쓴다} --- 위키
캐시가 유보 258 중 \textbf{6개}뿐이고 설명문도 6개다. 🔴 그 6은
\emph{``위키에 없다''} 가 아니라 \textbf{``안 봤다''} 다.

\begin{figure*}[t]
\centering
\includegraphics[width=\textwidth]{figs/fig1.pdf}
\caption{왼쪽: 자가 판에 닿는 데까지 200이 3으로 준다. 가운데: 팔 A --- 존재는
도메인마다 부호가 갈려 회색 띠(순열 널 $2\sigma$) 안에서 놀고, 양만 한쪽으로
선다. 오른쪽: 팔 B --- 신호 몫이 둘 다 문턱(붉은 파선)을 못 넘고, 잡음
바닥(점선)의 30\% 와 $-14$\% 자리에 있다.}
\label{fig:1}
\end{figure*}

\textbf{그리고 측정 전에 하나를 알았다.} \texttt{forms.py} 의 \texttt{\_feat} 는 축마다
\emph{값 열}과 \emph{관측 표시자 열}을 함께 낸다. \texttt{wikiaxes} 는 문서를 못 찾은
행의 마스크를 0 으로 준다. 그러면 \textbf{「위키 문서가 있다」는 이미 판에 표시자 열
셋으로 들어가 있다}.
\claimx{용량 ㄱ 의 개입 팔은 \textbf{새 정보 검정이 아니라 중복 검정}이다. 이것을
사후에 발견해 \emph{``모형이 이미 갖고 있다''} 고 주장하면 사후 선택이 된다 ---
그래서 사전등록에 예측 P3ㄱ 로 미리 박았다.}{위키 축을 base 에서 빼고 같은 존재 열을 주입했을 때 판이 오르면 '중복' 이 아니라 '표시자 열이 열등하다' 가 맞는 말이 된다 --- 이 사이클은 그걸 안 쟀다.}

\section{팔 A --- 존재는 죽고 양만 산다}

반응은 \textbf{행별 오차}다: 씨앗 0--5 순위 앙상블로
$e_i=|\mathrm{rank}(p_i)/n-\mathrm{rank}(y_i)/n|$. 통계는 도메인 안
$\rho_A=\mathrm{Spearman}(\text{용량},-e)$ 를 행수로 가중한 것이고, \textbf{널은 문턱
0.00353 이 아니라} 도메인 안 용량 라벨을 1,000회 순열해 다시 쟀다(그 팔의 $2\sigma$
는 존재 \textbf{0.0406} · 양 \textbf{0.0703} 으로 문턱의 11.5배 · 19.9배다 ---
P4 적중).

\begin{table}[h]\centering\small
\begin{tabular}{lrr}
\hline
 & ㄱ 존재 & ㄴ 양 \\
\hline
행 수(통계에 든 것) & 2,398 & 829 \\
$\rho_A$ 무통제 & $-0.0097$ & $\mathbf{+0.0692}$ \\
순열 널 $2\sigma$ & 0.0406 & 0.0703 \\
순열 널 양측 $p$ & 0.607 & \textbf{0.045} \\
BCa 95\% & $[-0.050,+0.033]$ & $\mathbf{[+0.002,+0.137]}$ \\
$+$글자수$+$어절수 & $-0.0126$ & $+0.0704$ \\
판정 & 판정 불능 & 승 \\
\hline
\end{tabular}
\end{table}

\claimx{존재는 도메인마다 부호가 갈린다 --- 도서 $+0.088$ · 모바일 $+0.062$ ·
게임 $+0.039$ · 세계애니 $+0.036$ 대 애니 $-0.082$ · 시장팝업 $-0.061$ ·
웹툰 $-0.048$. 모아서 $-0.0097$ 이고 순열 널 안이다.}{도메인을 늘리거나 만화·펀딩의 위키 덮음을 메운 뒤 다시 재서 $|\rho_A|$ 가 순열 널 $2\sigma$(0.0406) 밖으로 나오면 이 '판정 불능' 은 표본 부족이었다는 뜻이다.}
\claimx{양은 통제 뒤에도 산다 --- 글자수·어절수를 넣어도 $+0.0692\to+0.0704$
(\textbf{유지율 1.019}). P2(\emph{``통제하면 절반 이하로 준다''})는 \textbf{틀렸다}.
직전 노트의 실측(64--86\% 유지)이 맞았고 우리 예측이 틀렸다.}{인기·완결 여부·태그 수처럼 판이 이미 가진 열을 통제에 더 넣었을 때 $\rho_A$ 가 순열 널 안으로 들어가면, 이 상관은 텍스트가 아니라 레코드 충실도의 대용이었다.}

⚠ 팔 A 는 \textbf{인과가 아니다}. 그리고 원산국 통제는 원천이 세계애니뿐이라
\textbf{한 도메인}이다 --- \emph{``가중 편상관''} 이라 부르면 조항 60 위반이라
분모를 병기했다.

\section{팔 B --- 관찰에서 제일 센 도메인이 개입에서 제일 크게 잃는다}

개입은 표준 관례다: 도메인 안 백분위 1열 주입 · \textbf{값만 섞은 널 팔} 3뽑기 ·
씨앗 6 짝. 🔴 배선 검사를 \textbf{주입 뒤 \texttt{Data} 를 열어} 도메인마다 했다
(887 은 12도메인 중 아이돌만 조용히 중립화됐다). \textbf{887 형 중립화 0건}이고,
채택 안 한 도메인은 마스크 0 으로 \textbf{등록된 배제}다.

\begin{table}[h]\centering\small
\begin{tabular}{lrr}
\hline
 & ㄱ 존재 & ㄴ 양 \\
\hline
기준선 판 & 0.46917 & 0.46917 \\
진짜 판 & 0.46846 & 0.46837 \\
널 평균 & 0.46529 & 0.46989 \\
\textbf{신호 몫} & $+0.00318$ & $\mathbf{-0.00152}$ \\
\textbf{순효과} & $-0.00071$ & $-0.00080$ \\
씨앗 짝 BCa 95\% & $[+0.0015,+0.0052]$ & $[-0.0034,-0.0002]$ \\
문턱 0.00353 & \textbf{못 넘음}(0.90배) & 못 넘음 \\
잡음 바닥 0.01055 & 0.30배 & $-0.14$배 \\
\hline
\end{tabular}
\end{table}

\claimx{두 용량 다 \textbf{이 자를 못 넘었다}. 존재는 신호 몫이 문턱의 \textbf{0.90배}
이고 \textbf{순효과가 음수}다 --- 즉 진짜 열은 기준선을 낮추고, 섞은 열보다 덜 낮출
뿐이다. 양은 신호 몫이 \textbf{음수}이고 씨앗 짝 구간 상한이 0 아래다.}{널 뽑기를 셋에서 열둘로 늘리고 씨앗을 6에서 12로 늘렸을 때 존재의 신호 몫이 0.00353 위로 올라가면 이번 미달은 뽑기 잡음이었다.}
\claimx{🔴 \textbf{관찰과 개입이 정면으로 어긋난다.} 팔 A 에서 제일 센 도메인이
세계애니($\rho_A=+0.156$)인데 팔 B 에서 세계애니 순효과가 $\mathbf{-0.015}$ 로
제일 크게 손해다. 상관은 있는데 열로 주면 해롭다 --- 반증 조건 \textbf{F2 발화}:
모형이 이미 그 정보를 다른 열로 갖고 있다. \textbf{크롤은 중복이다.}}{설명문 \emph{본문}을 열로 넣었을 때(글자수가 아니라 내용) 세계애니 $\Delta\rho$ 가 양수로 뒤집히면, 중복인 것은 '텍스트' 가 아니라 '글자수' 다.}

⚠ 씨앗 짝 BCa 는 \textbf{씨앗} 재표집이라 행 잡음이 안 들어간다. 그리고 널 뽑기가
셋뿐이라 널 폭(존재 0.00204)이 신호 몫의 \textbf{64\%} 다 --- 구간을 액면 그대로
믿지 않는다.

\section{이 사이클이 정하는 것}

\claimx{\textbf{광역 크롤(12.5--34.4h, 댓글까지 229h)은 지금 근거가 없다.} 텍스트
\emph{용량}의 대용값 둘이 판을 문턱만큼 못 움직였고, 관찰에서 산 쪽은 개입에서
되레 해로웠다. 다만 이건 \emph{``텍스트를 넣으면 안 오른다''} 가 아니라
\emph{``텍스트가 있다/많다는 사실만으로는 안 오른다''} 다 --- \textbf{본문은 아직
한 번도 안 넣어 봤다}.}{파일럿 200건(0.3h)으로 본문을 긁어 열로 넣었을 때 만화 $\Delta\rho \ge 0.0517$ 이 나오면 크롤에 근거가 생긴다 --- 그 수는 이 노트가 0단계에서 계산해 뒀다.}
\claimx{다음 자는 \textbf{유보에서 표본을 뽑아야} 하고, 만화의 위키·설명문
덮음(258 중 6)을 먼저 메워야 어떤 자도 만화를 못 잰다.}{만화 유보 258 에 위키·설명문을 채운 뒤에도 $\rho_A$ 가 순열 널 안이면, 만화가 안 잡히는 이유는 덮음이 아니라 도메인 자체다.}

\section{자가 적발}

\claimx{🔴 초판 팔 A 가 \textbf{게임과 웹툰을 조용히 중립화}했다. \texttt{scipy
1.13} 의 \texttt{rankdata} 는 배열에 NaN 이 \textbf{하나라도} 있으면 \textbf{전부}
NaN 을 낸다 --- 라벨 결측이 게임 43건 · 웹툰 61건 있어서 두 도메인의 오차가 통째로
NaN 이 됐고, \emph{``채택 도메인''} 목록에 이름은 남은 채 통계에는 한 행도 안
들어갔다. 887 형 중립화의 통계 판본이다. 유한한 행 안에서 순위를 다시 내도록 고쳤고
$\rho_A$ 가 $+0.0006\to-0.0097$ 로 바뀌었다(판정은 불변 --- 둘 다 널 안).
\textbf{회계 칸(``통계에 실제로 든 행'')을 산출물에 새로 박았다.}}{저장소에서 \texttt{rankdata} 를 거르지 않고 부르는 다른 자리를 세어 0 이 아니면 같은 병이 더 있다 --- 이 노트는 그것을 안 셌다.}

\end{document}
